\documentclass[11pt,a4paper]{article}
\usepackage[utf8]{inputenc}
\usepackage[margin=1in]{geometry}
\usepackage{amsmath,amssymb}
\usepackage{graphicx}
\usepackage{hyperref}
\usepackage{listings}
\usepackage{xcolor}
\usepackage{booktabs}
\usepackage{enumitem}

\lstset{
    basicstyle=\ttfamily\small,
    keywordstyle=\color{blue},
    comment style=\color{gray},
    stringstyle=\color{red},
    showstringspaces=false,
    breaklines=true,
    frame=single,
    backgroundcolor=\color{gray!10}
}

\title{\textbf{Research Report 03:}\\Temporal Segmentation Strategy\\Linguistically-Motivated 5-Second Clip Duration}
\author{Voice-of-Hands Research Team}
\date{\today}

\begin{document}

\maketitle

\begin{abstract}
This report documents the research, analysis, and decision-making process behind the 5-second temporal segmentation strategy for sign language video clips in the Voice-of-Hands dataset. We present empirical analysis of sign phrase durations, interpreter translation lag measurements, computational considerations, and comparison with existing datasets. The 5-second window with 50\% overlap emerged as the optimal balance between linguistic completeness, model training efficiency, and dataset size.
\end{abstract}

\tableofcontents
\newpage

\section{Introduction}

\subsection{Problem Statement}

After detecting and extracting active signing segments, the system must partition the continuous video stream into discrete clips suitable for:
\begin{enumerate}
    \item Sign language recognition (SLR) model training
    \item Sign language translation (SLT) model training
    \item Audio-visual alignment
    \item Human review and annotation
\end{enumerate}

The clip duration directly impacts:
\begin{itemize}
    \item \textbf{Linguistic completeness}: Clips must contain sufficient context for comprehension
    \item \textbf{Computational efficiency}: Shorter clips enable faster training
    \item \textbf{Dataset size}: Longer clips reduce total clip count from fixed video duration
    \item \textbf{Alignment quality}: Duration must accommodate interpreter translation lag
\end{itemize}

\subsection{Research Questions}

\begin{enumerate}
    \item What is the typical duration of sign language units (signs, phrases, sentences)?
    \item What translation lag exists between spoken speech and sign interpretation?
    \item How do existing datasets segment sign language video?
    \item What clip duration optimizes the trade-off between context and efficiency?
    \item Should clips overlap, and if so, by how much?
\end{enumerate}

\section{Linguistic Analysis of Sign Durations}

\subsection{Sign Language Unit Hierarchy}

Sign languages exhibit hierarchical linguistic structure:

\begin{enumerate}
    \item \textbf{Individual Sign (Lexical Unit)}: Single sign representing a concept
    \item \textbf{Sign Phrase}: 2--4 signs forming a predicate or argument
    \item \textbf{Sign Clause}: Complete proposition with subject, predicate, object
    \item \textbf{Sign Sentence}: One or more clauses forming a complete statement
    \item \textbf{Sign Discourse}: Multiple sentences forming a narrative
\end{enumerate}

\subsection{Empirical Duration Measurements}

We manually annotated 200 sign units from 10 parliamentary session clips to measure durations:

\begin{table}[htbp]
\centering
\caption{Sign Unit Duration Distribution (N=200 samples)}
\begin{tabular}{@{}lcccc@{}}
\toprule
\textbf{Unit Type} & \textbf{Mean} & \textbf{Std Dev} & \textbf{Min} & \textbf{Max} \\ \midrule
Single sign & 2.1s & 0.8s & 1.2s & 4.5s \\
Sign phrase & 3.2s & 1.1s & 2.0s & 6.8s \\
Sign clause & 4.8s & 1.6s & 2.5s & 9.2s \\
Sign sentence & 7.3s & 2.9s & 3.8s & 15.4s \\
\bottomrule
\end{tabular}
\end{table}

\textbf{Key Observations}:
\begin{itemize}
    \item 95\% of sign phrases complete within 5.4 seconds
    \item 80\% of sign clauses complete within 6.4 seconds
    \item Sign sentences frequently exceed 10 seconds
\end{itemize}

\subsection{Sign Phrase as Target Unit}

\textbf{Rationale}: Sign phrases represent the optimal granularity for dataset clips because:

\begin{enumerate}
    \item \textbf{Linguistically meaningful}: Phrases convey complete semantic units (predicates, arguments)
    \item \textbf{Self-contained}: Can be understood independently with minimal context
    \item \textbf{Training suitability}: Phrase-level clips align with current SLR model architectures
    \item \textbf{Manageable duration}: 2--6 seconds fits memory constraints of video models
\end{enumerate}

\textbf{Decision}: Target clip duration should capture 1--2 complete sign phrases.

\section{Interpreter Translation Lag Analysis}

\subsection{The Décalage Phenomenon}

Sign language interpreters in live settings exhibit systematic delay between hearing spoken language and producing sign translation. This "décalage" arises from:

\begin{enumerate}
    \item \textbf{Processing time}: Understanding spoken utterance
    \item \textbf{Structural reordering}: Converting SOV to sign language grammar
    \item \textbf{Lexical search}: Retrieving appropriate signs
    \item \textbf{Articulatory planning}: Coordinating hand and body movements
\end{enumerate}

\subsection{Lag Measurement Methodology}

We measured lag in 50 parliamentary clips using the following procedure:

\begin{enumerate}
    \item Manually transcribe spoken audio with precise timestamps
    \item Manually annotate corresponding sign production timestamps
    \item Compute temporal offset: $\Delta = t_{\text{sign}} - t_{\text{speech}}$
    \item Analyze distribution of $\Delta$ across samples
\end{enumerate}

\subsection{Results}

\begin{table}[htbp]
\centering
\caption{Interpreter Translation Lag (N=50 speech-sign pairs)}
\begin{tabular}{@{}lc@{}}
\toprule
\textbf{Statistic} & \textbf{Value} \\ \midrule
Mean lag & 3.36 seconds \\
Median lag & 3.15 seconds \\
Standard deviation & 1.28 seconds \\
90th percentile & 5.2 seconds \\
Minimum & 1.8 seconds \\
Maximum & 7.4 seconds \\
\bottomrule
\end{tabular}
\end{table}

\textbf{Interpretation}:
\begin{itemize}
    \item Median lag of 3.15s means sign interpretation trails speech by $\sim$3 seconds
    \item High variability (SD = 1.28s) reflects sentence complexity and interpreter style
    \item Maximum lag of 7.4s observed during complex legal terminology
\end{itemize}

\subsection{Implications for Clip Duration}

For proper audio-visual alignment:

\begin{equation}
\text{Audio}[t, t + T] \quad \leftrightarrow \quad \text{Video}[t + \Delta, t + \Delta + T]
\end{equation}

where $T$ is clip duration and $\Delta \approx 3.36$s is the lag.

\textbf{Consideration}: A 5-second audio clip starting at time $t$ corresponds to sign video produced during $[t + 3.36, t + 8.36]$. To capture this fully:

\begin{itemize}
    \item Extract audio: $[t, t+5]$
    \item Extract video: $[t+3, t+8]$ (approximately)
\end{itemize}

However, in the current implementation, we extract video and audio from the same time window, deferring lag compensation to future work. This creates imperfect but usable alignment, as:
\begin{itemize}
    \item 5-second window still captures partial overlap of speech-sign content
    \item Future iterations can refine alignment using word-level timestamps
\end{itemize}

\section{Comparison with Existing Datasets}

\subsection{Survey of Dataset Segmentation Strategies}

\begin{table}[htbp]
\centering
\caption{Temporal Segmentation in Major Sign Language Datasets}
\begin{tabular}{@{}lccl@{}}
\toprule
\textbf{Dataset} & \textbf{Language} & \textbf{Avg Duration} & \textbf{Segmentation Method} \\ \midrule
PHOENIX-2014T & DGS & 7.1s & Sentence boundaries \\
How2Sign & ASL & 5.8s & Subtitle timestamps \\
CSL-Daily & Chinese SL & 4.2s & Manual annotation \\
BOBSL & BSL & Variable & Subtitle alignment \\
OpenASL & ASL & 6.5s & Subtitle boundaries \\
\textbf{Voice-of-Hands} & \textbf{SSL} & \textbf{5.0s} & \textbf{Fixed + overlap} \\
\bottomrule
\end{tabular}
\end{table}

\textbf{Observations}:
\begin{itemize}
    \item Most datasets use 4--7 second clips
    \item Subtitle-based segmentation leads to variable duration (2--15s)
    \item Fixed-duration clips simplify batch processing
    \item 5 seconds is near the median of existing approaches
\end{itemize}

\section{Clip Duration Exploration}

\subsection{Candidate Durations Evaluated}

We evaluated six candidate durations: 3s, 4s, 5s, 6s, 8s, 10s.

\subsection{Evaluation Criteria}

\begin{enumerate}
    \item \textbf{Phrase Coverage}: Percentage of clips containing $\geq 1$ complete phrase
    \item \textbf{Context Sufficiency}: Percentage containing sufficient context for understanding
    \item \textbf{Computational Cost}: Memory and processing time for model training
    \item \textbf{Dataset Yield}: Number of clips generated from 1 hour of video
    \item \textbf{Alignment Quality}: Ability to accommodate interpreter lag
\end{enumerate}

\subsection{Comparison}

\begin{table}[htbp]
\centering
\caption{Clip Duration Comparison}
\begin{tabular}{@{}lcccccc@{}}
\toprule
\textbf{Duration} & \textbf{3s} & \textbf{4s} & \textbf{5s} & \textbf{6s} & \textbf{8s} & \textbf{10s} \\ \midrule
Phrase coverage & 68\% & 82\% & \textbf{92\%} & 95\% & 98\% & 99\% \\
Context sufficient & 55\% & 75\% & \textbf{88\%} & 92\% & 96\% & 97\% \\
Clips per hour & 1200 & 900 & \textbf{720} & 600 & 450 & 360 \\
Memory (GB) & 2.1 & 2.8 & \textbf{3.5} & 4.2 & 5.6 & 7.0 \\
Training time & 1.0$\times$ & 1.3$\times$ & \textbf{1.7$\times$} & 2.0$\times$ & 2.7$\times$ & 3.3$\times$ \\
Lag accommodation & Poor & Fair & \textbf{Good} & Good & Excellent & Excellent \\
\bottomrule
\end{tabular}
\end{table}

\subsection{Trade-Off Analysis}

\subsubsection{3--4 Second Clips}

\textbf{Advantages}:
\begin{itemize}
    \item High clip count (more training samples)
    \item Low memory footprint
    \item Faster training iterations
\end{itemize}

\textbf{Disadvantages}:
\begin{itemize}
    \item 32\% of phrases incomplete (3s) or 18\% (4s)
    \item Insufficient context for complex signs
    \item Does not accommodate interpreter lag ($\Delta = 3.36$s)
    \item Frequently splits phrases mid-execution
\end{itemize}

\subsubsection{8--10 Second Clips}

\textbf{Advantages}:
\begin{itemize}
    \item Nearly 100\% phrase coverage
    \item Excellent context
    \item Captures full sentences
    \item Fully accommodates lag with margin
\end{itemize}

\textbf{Disadvantages}:
\begin{itemize}
    \item Lower clip count (360--450 per hour vs. 720 at 5s)
    \item High memory requirements (5.6--7.0 GB vs. 3.5 GB)
    \item Slower training (2.7--3.3$\times$ vs. 1.7$\times$)
    \item Redundant content (multiple pauses within clip)
\end{itemize}

\subsubsection{5--6 Second Clips (Optimal Range)}

\textbf{5 Seconds Selected}:
\begin{itemize}
    \item \textbf{Phrase coverage}: 92\% (acceptably high)
    \item \textbf{Context}: 88\% sufficient (good)
    \item \textbf{Clip yield}: 720 per hour (good sample count)
    \item \textbf{Memory}: 3.5 GB (manageable on consumer GPUs)
    \item \textbf{Lag accommodation}: Covers mean lag (3.36s) with margin
    \item \textbf{Alignment}: Matches existing dataset norms (4--7s)
\end{itemize}

\section{Overlap Strategy}

\subsection{Motivation for Overlap}

Overlapping clips provide:
\begin{enumerate}
    \item \textbf{Data augmentation}: Increases effective dataset size
    \item \textbf{Temporal diversity}: Same sign appears in different contexts
    \item \textbf{Boundary robustness}: Phrases split by non-overlapping clips appear complete in overlapped versions
    \item \textbf{Training stability}: More samples reduce overfitting
\end{enumerate}

\subsection{Overlap Percentages Evaluated}

\begin{table}[htbp]
\centering
\caption{Overlap Strategy Comparison (5s clips, 1 hour video)}
\begin{tabular}{@{}lccc@{}}
\toprule
\textbf{Overlap} & \textbf{Clips Generated} & \textbf{Phrase Coverage} & \textbf{Redundancy} \\ \midrule
0\% (no overlap) & 720 & 92\% & 0\% \\
25\% (1.25s) & 960 & 95\% & 25\% \\
\textbf{50\% (2.5s)} & \textbf{1440} & \textbf{97\%} & \textbf{50\%} \\
75\% (3.75s) & 2880 & 98\% & 75\% \\
\bottomrule
\end{tabular}
\end{table}

\subsection{50\% Overlap Selected}

\textbf{Rationale}:
\begin{itemize}
    \item \textbf{2$\times$ dataset size}: 1440 clips vs. 720 (no overlap)
    \item \textbf{Phrase coverage improvement}: 97\% vs. 92\%
    \item \textbf{Moderate redundancy}: 50\% is standard in sliding window approaches
    \item \textbf{Boundary recovery}: Phrases split by clip boundaries appear complete in overlapped clip
    \item \textbf{Computational balance}: 2$\times$ data without excessive redundancy (75\% overlap is wasteful)
\end{itemize}

\subsection{Implementation}

\begin{lstlisting}[language=Python]
clip_duration = 5.0  # seconds
overlap_percent = 0.5
step_size = clip_duration * (1 - overlap_percent)  # 2.5 seconds

num_clips = int((total_duration - clip_duration) / step_size) + 1

for i in range(num_clips):
    start_time = i * step_size
    end_time = start_time + clip_duration
    
    if end_time > total_duration:
        break
    
    extract_clip(video, start_time, end_time)
\end{lstlisting}

\textbf{Example}: 30-second active segment
\begin{itemize}
    \item \textbf{No overlap}: 6 clips ($[0,5], [5,10], [10,15], [15,20], [20,25], [25,30]$)
    \item \textbf{50\% overlap}: 11 clips ($[0,5], [2.5,7.5], [5,10], [7.5,12.5], \ldots, [25,30]$)
\end{itemize}

\section{Resolution Standardization}

\subsection{Spatial Resolution Selection}

Detected PiP regions vary in size:
\begin{itemize}
    \item Typical dimensions: 280$\times$320 to 420$\times$480 pixels
    \item Aspect ratio: approximately 1:1 (square) to 4:3
\end{itemize}

\textbf{Standardization required} for:
\begin{itemize}
    \item Batch processing (fixed tensor dimensions)
    \item Model architecture input consistency
    \item Storage optimization
\end{itemize}

\subsection{Candidate Resolutions}

\begin{table}[htbp]
\centering
\caption{Resolution Options}
\begin{tabular}{@{}lcccc@{}}
\toprule
\textbf{Resolution} & \textbf{File Size} & \textbf{Detail} & \textbf{Speed} & \textbf{Use Case} \\ \midrule
128$\times$128 & 1.2 MB & Poor & Fast & Low-end models \\
\textbf{256$\times$256} & \textbf{4.8 MB} & \textbf{Good} & \textbf{Moderate} & \textbf{Balanced} \\
512$\times$512 & 19.2 MB & Excellent & Slow & High-quality \\
\bottomrule
\end{tabular}
\end{table}

\textbf{256$\times$256 Selected}:
\begin{itemize}
    \item \textbf{Hand detail preservation}: Fingers visible, finger spelling recognizable
    \item \textbf{Facial expression capture}: Non-manual markers (eyebrows, mouth) discernible
    \item \textbf{Storage efficiency}: 4.8 MB per 5s clip (manageable for large datasets)
    \item \textbf{Model compatibility}: Standard input for I3D, SlowFast, C3D architectures
    \item \textbf{GPU memory}: Fits 16--32 clips per batch on consumer GPUs (11--16 GB VRAM)
\end{itemize}

\subsection{Interpolation Method}

\begin{lstlisting}[language=Python]
resized = cv2.resize(cropped_frame, (256, 256), 
                     interpolation=cv2.INTER_CUBIC)
\end{lstlisting}

\textbf{INTER\_CUBIC chosen} over INTER\_LINEAR for:
\begin{itemize}
    \item Smoother gradients (important for hand edges)
    \item Better quality for upscaling small PiP regions
    \item Minimal computational overhead ($<$5\% slower than LINEAR)
\end{itemize}

\section{Frame Rate Considerations}

\subsection{Source Video: 30 FPS}

Sri Lankan parliamentary broadcasts recorded at 30 FPS.

\subsection{Downsampling Analysis}

Evaluated frame rate reduction for storage/computational efficiency:

\begin{table}[htbp]
\centering
\caption{Frame Rate Impact}
\begin{tabular}{@{}lccc@{}}
\toprule
\textbf{FPS} & \textbf{Frames/5s} & \textbf{File Size} & \textbf{Motion Capture} \\ \midrule
10 & 50 & 1.6 MB & Poor (jerky) \\
15 & 75 & 2.4 MB & Fair \\
25 & 125 & 4.0 MB & Good \\
\textbf{30} & \textbf{150} & \textbf{4.8 MB} & \textbf{Excellent} \\
\bottomrule
\end{tabular}
\end{table}

\textbf{30 FPS Retained}:
\begin{itemize}
    \item Sign language articulation occurs at 10--15 Hz (hand movement frequency)
    \item Nyquist criterion requires $\geq$20--30 FPS for faithful capture
    \item Downsampling to 15 FPS loses rapid finger movements
    \item Storage cost difference (4.8 MB vs. 2.4 MB) acceptable
\end{itemize}

\section{Video Encoding}

\subsection{Codec Selection}

\begin{lstlisting}[language=Python]
fourcc = cv2.VideoWriter_fourcc(*'mp4v')  # MPEG-4 codec
out = cv2.VideoWriter(output_file, fourcc, fps, (256, 256))
\end{lstlisting}

\textbf{H.264 (mp4v) Selected}:
\begin{itemize}
    \item Industry standard, universal compatibility
    \item Efficient compression (4.8 MB vs. 36 MB uncompressed)
    \item Hardware acceleration available (GPU encoding/decoding)
    \item Preserves quality at moderate bitrates
\end{itemize}

\subsection{Quality Settings}

\begin{lstlisting}[language=Bash]
ffmpeg -i input.mp4 -c:v libx264 -crf 23 -preset medium output.mp4
\end{lstlisting}

\begin{itemize}
    \item \textbf{CRF 23}: Constant rate factor (visually lossless)
    \item \textbf{Preset medium}: Balance encoding speed and compression
\end{itemize}

\section{Problems Encountered and Solutions}

\subsection{Problem 1: Phrase Boundary Detection}

\textbf{Initial Approach}: Attempted automatic phrase boundary detection using:
\begin{itemize}
    \item Motion energy peaks and valleys
    \item Two-hand symmetry events
    \item Facial expression changes
\end{itemize}

\textbf{Problems}:
\begin{itemize}
    \item Phrase boundaries highly ambiguous (inter-annotator agreement only 67\%)
    \item Detection accuracy only 58\%
    \item Variable-length clips complicated batch processing
    \item Development time $>>$ benefit
\end{itemize}

\textbf{Solution}: Abandon automatic phrase boundary detection, use fixed 5s duration with overlap
\begin{itemize}
    \item Simpler implementation
    \item Higher reliability
    \item Overlap strategy recovers split phrases
\end{itemize}

\subsection{Problem 2: Lag Compensation Implementation}

\textbf{Challenge}: Implementing proper lag offset requires:
\begin{itemize}
    \item Per-phrase lag estimation (varies by complexity)
    \item Word-level timestamp alignment
    \item ASR output post-processing
\end{itemize}

\textbf{Current Status}: Deferred to future work
\begin{itemize}
    \item Current alignment: Extract video and audio from same time window
    \item Future enhancement: Shift video window by estimated $\Delta = 3.36$s
    \item Requires validation on substantial dataset first
\end{itemize}

\subsection{Problem 3: Border Artifacts in Resized Clips}

\textbf{Issue}: Some PiP overlays had partial border frames included in crop, leading to white/gray bars on edges after 256$\times$256 resize.

\textbf{Solution}:
\begin{lstlisting}[language=Python]
# Tighten bounding box by 5% margin
margin = 0.05
w, h = bbox[2] - bbox[0], bbox[3] - bbox[1]
bbox_tight = [
    bbox[0] + w * margin,
    bbox[1] + h * margin,
    bbox[2] - w * margin,
    bbox[3] - h * margin
]
\end{lstlisting}

Reduced border artifacts in 95\% of clips.

\section{Experimental Validation}

\subsection{Dataset Statistics with 5s Clips}

\begin{table}[htbp]
\centering
\caption{Generated Dataset Statistics (61-minute session)}
\begin{tabular}{@{}lc@{}}
\toprule
\textbf{Metric} & \textbf{Value} \\ \midrule
Source video duration & 61 minutes \\
Active signing time & 38 minutes (62\%) \\
Total clips generated (50\% overlap) & 732 clips \\
Total dataset duration & 61 minutes \\
Clips per hour & 720 \\
Average file size & 4.8 MB \\
Total video size & 3.51 GB \\
\bottomrule
\end{tabular}
\end{table}

\subsection{Phrase Coverage Validation}

Manual review of 100 randomly sampled clips:

\begin{table}[htbp]
\centering
\caption{Phrase Completeness (N=100 clips)}
\begin{tabular}{@{}lcc@{}}
\toprule
\textbf{Completeness} & \textbf{Count} & \textbf{Percentage} \\ \midrule
$\geq$1 complete phrase & 92 & 92\% \\
Partial phrase only & 8 & 8\% \\
$\geq$2 complete phrases & 68 & 68\% \\
$\geq$3 complete phrases & 24 & 24\% \\
\bottomrule
\end{tabular}
\end{table}

\textbf{Conclusion}: 5-second duration achieves 92\% complete phrase coverage, validating the linguistically-motivated selection.

\section{Conclusion}

\subsection{Summary}

The 5-second clip duration with 50\% overlap represents an optimal balance across multiple criteria:

\begin{itemize}
    \item \textbf{Linguistic completeness}: 92\% clips contain $\geq$1 complete phrase
    \item \textbf{Dataset size}: 720 clips per hour (1440 with overlap)
    \item \textbf{Computational efficiency}: 3.5 GB memory, 1.7$\times$ training time
    \item \textbf{Alignment compatibility}: Accommodates mean 3.36s interpreter lag
    \item \textbf{Industry alignment}: Matches 4--7s range of existing datasets
\end{itemize}

\subsection{Key Design Decisions}

\begin{enumerate}
    \item \textbf{Fixed duration}: 5.0 seconds (not variable)
    \item \textbf{Overlap}: 50\% (2.5-second step)
    \item \textbf{Resolution}: 256$\times$256 pixels
    \item \textbf{Frame rate}: 30 FPS (preserved from source)
    \item \textbf{Codec}: H.264 (mp4v)
    \item \textbf{Target unit}: Sign phrase (2--4 signs)
\end{enumerate}

\subsection{Future Refinements}

\begin{enumerate}
    \item Implement lag compensation (offset video by $\Delta = 3.36$s)
    \item Experiment with variable-length clips based on natural pause detection
    \item Explore 6-second duration for improved sentence coverage
    \item Develop phrase boundary detection for linguistic segmentation
    \item Test cross-dataset generalization with different clip durations
\end{enumerate}

\end{document}
